\documentclass[11pt]{article}
\usepackage[margin=1in]{geometry}
\usepackage{amsmath,amssymb,amsthm,booktabs,microtype}
\usepackage[colorlinks=true,urlcolor=blue,citecolor=blue]{hyperref}
\newtheorem{theorem}{Theorem}
\newtheorem{proposition}{Proposition}
\newtheorem{lemma}{Lemma}
\newcommand{\rhoo}{\rho}
\title{Three-way assignment can separate fair core selections maximally}
\author{Internal research note}
\date{4 September 2026}
\begin{document}
\maketitle
\begin{abstract}
We construct standard three-sided assignment games with nonempty entire cores,
equally sized sectors, and triple values in $\{0,1,2,3\}$ whose full leximin and
full descending-lexicographic-minimum payoffs have normalized $\ell_1$ distance
tending to two. This is the sharp supremum. The same construction rules out
every positive dimension-independent lower guarantee for a player's leximin
payoff relative to their maximum core payoff. A companion family has
exponentially vanishing individual guarantees but coincident selectors.
The proof uses explicit integral matchings and the entire core, with every
fixed-zero player counted. Correctness has been checked internally;
novelty and external significance remain subject to specialist review.
\end{abstract}

\section{Statement and model}
Let $A,B,C$ be finite disjoint sets and $w_{abc}\ge0$. A coalition's worth is
the maximum weight of an \emph{integral} matching of disjoint triples within
it. Write $T>0$ for the grand coalition's matching value. Its entire core is
\[
 K=\{p\ge0:\ \sum_i p_i=T,\quad p_a+p_b+p_c\ge w_{abc}
       \text{ for every }(a,b,c)\}.
\]
Indeed, the core implies each triple inequality, while summing these inequalities
over a coalition's optimal matching proves every coalition inequality.
We assume $K\ne\varnothing$; nonemptiness does not hold for arbitrary arrays.
Let $x$ maximize the ascending sorted payoff vector lexicographically, and let
$z$ minimize the descending sorted vector lexicographically. Both are full
selectors; fixed zeros remain in the sorted vectors. Set
$\rho=\|x-z\|_1/T$.

\begin{theorem}\label{thm:main}
For every integer $L\ge1$ and $R\ge0$, there is such a game with
\[
 n=9(L+R)+12,\qquad T=3(L+R+2),\qquad w_{abc}\in\{0,1,2,3\},
\]
and
\begin{equation}\label{eq:rho}
 \rho(L,R)=\frac{2((R+2)2^L-1)}{(2^L+2)(L+R+2)}.
\end{equation}
Consequently the sharp supremum of $\rho$ over three-sided assignment games
with nonempty entire cores is $2$. A player in this family has maximum core
payoff $1$ but leximin payoff $2/(2^L+2)$, so no positive dimension-independent
individual fraction survives.
\end{theorem}

The two-sided predecessor establishes $x_i\ge y_i/2$ for every core point $y$
and sharp maximal distance $\lfloor n/2\rfloor/(\lfloor n/2\rfloor+1)$ for
$n\ge4$; that internal result is preserved separately. Theorem~\ref{thm:main}
places the change at team size three, even with bounded integer values.
It neither asserts a sharp finite-$n$ formula nor resolves the best convergence rate.

\section{An explicit market}
All color indices are modulo three. Introduce global players $G_i$ of color
$i$. At levels $\ell=0,\ldots,L$ introduce $U_{\ell i},V_{\ell i},D_{\ell i}$
for $i=0,1,2$. At even levels their colors are $i,i+1,i+2$; at odd levels
the colors of $U$ and $V$ are interchanged. Put $s_0=2$ and $s_\ell=1$
for $\ell\ge1$. Assign the following positive triple values:
\begin{center}
\begin{tabular}{ll}
\toprule
Triple & Value\\
\midrule
$\{U_{\ell i},V_{\ell i},D_{\ell i}\}$ & $s_\ell$\\
$\{U_{\ell i},V_{\ell i},G_{i+2}\}$ & $s_\ell$\\
$\{U_{0i},G_{i+1},G_{i+2}\}$ & $1$\\
$\{U_{\ell i},V_{\ell-1,j},V_{\ell-1,k}\}$, $\ell\ge1$ & $2$\\
$\{V_{\ell i},U_{\ell-1,j},U_{\ell-1,k}\}$, $\ell\ge1$ & $3$ if $\ell=1$, else $1$\\
\bottomrule
\end{tabular}
\end{center}
In the last two rows, $\{j,k\}=\{0,1,2\}\setminus\{i\}$.
Each listed triple has exactly one player of each color.

For every $i$ introduce $R$ copies $(U'_{ri},V'_{ri},D'_{ri})$ with the same
colors as its top-level triple. Give value $1$ to each copy's local and global
pair triples $\{U'_{ri},V'_{ri},D'_{ri}\}$ and
$\{U'_{ri},V'_{ri},G_{i+2}\}$, and to its cross-triples
\[
 \{U'_{ri},V_{Li},G_{i+2}\},\qquad
 \{V'_{ri},U_{Li},G_{i+2}\}.
\]
Every unlisted triple has value zero. There are $3(L+R)+4$ players per sector.

\section{Integral value and entire core}
The triples using each pair's own $D$ form a disjoint integral matching with
weight $6+3L+3R=T$. Consider the allocation assigning $1$ to both root
$U$ and root $V$, $0$ to every subsequent $U$, $1$ to every subsequent $V$,
the analogous values to copies, and $0$ to every dummy. It has total $T$
and covers every weighted triple. Summing its payoffs over any disjoint
matching bounds that matching's weight by $T$. Hence the displayed integral
matching is optimal and this allocation belongs to the entire core.

\begin{lemma}\label{lem:core}
Every core allocation, and only a core allocation, has all dummy payoffs zero
and the following form:
\begin{align*}
 U_{0i}&=1+q_{0i}, & V_{0i}&=1-q_{0i},\\
 U_{\ell i}&=q_{\ell i}, &V_{\ell i}&=1-q_{\ell i}\quad(\ell\ge1),\\
 U'_{ri}&=q_{Li}, &V'_{ri}&=1-q_{Li},
\end{align*}
where $0\le q_{\ell i}\le1$ and
\begin{equation}\label{eq:rec}
 q_{\ell i}=q_{\ell-1,j}+q_{\ell-1,k},\qquad
 q_\ell=(J-I)q_{\ell-1}.
\end{equation}
Here $J$ is the $3\times3$ all-ones matrix.
\end{lemma}
\begin{proof}
Sum the inequalities of the displayed matching. Efficiency and nonnegativity
force every unmatched $G$ to zero and every matching triple to equality.
The alternative pair triple using $G$ then forces its local $D$ to zero and
its pair to total $s_\ell$. Root lower-bound triples give $U_{0i}\ge1$.
At level one, the two recurrence inequalities become respectively
$q_{1i}\ge q_{0j}+q_{0k}$ and $q_{1i}\le q_{0j}+q_{0k}$ after substituting
the root offsets. At subsequent levels the same substitutions give these
opposite inequalities without offsets. Cross-triples give opposite inequalities
between a copied $U$ and its original top $U$, hence equality.
Nonnegativity supplies the stated bounds.

Conversely, the parameterization satisfies every displayed triple inequality,
has total $T$, and is nonnegative. All omitted zero-weight triple inequalities
are automatic. It therefore satisfies every coalition inequality by summing
over a matching. This proves equality with the entire core.
\end{proof}

\section{The full selectors and quantitative separation}
For minimax, every root $U$ is at least $1$. The allocation just used, denoted
$z$, has maximum payoff $1$. Any core allocation with maximum at most $1$
must have $U_{0i}=1$ for all $i$, hence $q_0=0$. Lemma~\ref{lem:core} then
fixes every coordinate. Thus $z$ is the unique minimax allocation and the
full descending-lexicographic minimum.

For leximin, all $D,G$ are fixed at zero. Every other coordinate can
simultaneously be positive by taking $q_0=t(1,1,1)$ with $0<t<2^{-L}$.
Consequently the full procedure first fixes exactly the dummy coordinates
at zero, then maximizes a common floor $\eta$ for all remaining coordinates.
If such a floor is feasible, level-one $U$ and top-level $V$ imply
\[
 2\sum_iq_{0i}=\sum_iq_{1i}\ge3\eta,\qquad
 2^L\sum_iq_{0i}=\sum_iq_{Li}\le3(1-\eta).
\]
Combining gives $\eta\le2/(2^L+2)$. Equality is attained by
\begin{equation}\label{eq:x}
 t=\frac1{2^L+2},\qquad q_{\ell i}=2^\ell t.
\end{equation}
The least nonroot $U$ is $2t$ and the least $V$ is $1-2^Lt=2t$;
root and copy payoffs also respect this floor. Equality forces
$q_{1i}=\eta$ for all three $i$. Since $J-I$ is invertible, with eigenvalues
$2,-1,-1$, it then forces $q_0=t(1,1,1)$. Thus the optimal positive-floor
face is a singleton and \eqref{eq:x} gives the full leximin allocation $x$.
No symmetry restriction on the feasible core was imposed.

Every level-$\ell$ pair changes from $z$ to $x$ by
$(2^\ell t,-2^\ell t)$, including the offset root pair. Each copied pair
changes by $(2^Lt,-2^Lt)$. Therefore
\[
 \|x-z\|_1=6t\left(\sum_{\ell=0}^{L}2^\ell+R2^L\right)
 =\frac{6((R+2)2^L-1)}{2^L+2}.
\]
Dividing by $T$ proves \eqref{eq:rho}. With $R=L^2$, both the amplified
payoff fraction and the share of copied pairs tend to one, giving $\rho\to2$.
Conversely, nonnegative vectors of total $T$ always satisfy
$\|x-z\|_1\le2T$, proving the sharp supremum.
For example, $L=4,R=16$ gives $\rho=287/198>4/3$.
Finally, every top $V$ has core maximum $1$, attained at $z$, but payoff
$2/(2^L+2)$ at $x$. This proves the individual conclusion separately.

\begin{proposition}\label{prop:dimension}
Let $S_n$ be the supremum for total player count $n$, allowing arbitrary sector
sizes. For $n\ge30$, put $k=\lfloor(n-12)/9\rfloor$. Then
\[
 0\le2-S_n\le
 \frac{2\lceil\log_2 k\rceil+4}{k+2}+\frac{2}{(k+2)^2}
 =O(\log n/n).
\]
\end{proposition}
\begin{proof}
Choose $L=\lceil\log_2 k\rceil$ and $R=k-L\ge0$, and add the remaining
at most eight players with zero incident values. Matching efficiency forces
these players to zero in the entire core, preserving both full selectors.
Subtracting \eqref{eq:rho} from two gives
\[
 2-\rho=\frac{2L2^L}{(2^L+2)(k+2)}
       +\frac4{2^L+2}+\frac2{(2^L+2)(k+2)}.
\]
Use $2^L\ge k$ and bound the first fraction by $2L/(k+2)$.
\end{proof}

\section{Controls, calibration, and scope}
The construction separates two mechanisms. Recurrence amplifies payoff
differences across levels. Root pair offsets pin the minimax allocation;
top copies make the amplified changes dominate total value. Taking $R=0$
and $L\to\infty$ makes both the individual ratio and $\rho$ tend to zero.
Taking fixed $L$ and $R\to\infty$ gives only
$\rho\to2^{L+1}/(2^L+2)<2$.

There is also an exact coincident-selector control. Set every pair total to
$1$, remove root lower-bound triples and root offsets, and change the
level-one complementary triple value from $3$ to $1$. With no copies the
same proof gives $U_\ell=(J-I)^\ell U_0$, $V_\ell=1-U_\ell$.
A positive floor $\eta$ gives $\sum_iU_{0i}\ge3\eta$ and
$2^L\sum_iU_{0i}\le3(1-\eta)$, hence $\eta\le1/(2^L+1)$.
A ceiling $M$ gives $\sum_iU_{0i}\ge3(1-M)$ and
$2^L\sum_iU_{0i}\le3M$, hence $M\ge2^L/(2^L+1)$.
Both bounds are attained by $U_{\ell i}=2^\ell/(2^L+1)$; equality in
each pins all root coordinates and therefore the entire allocation.
Thus $x=z$, while a top $V$ has selected payoff
$1/(2^L+1)$ and core maximum $1$. The proposed one-third individual
guarantee fails already at $L=2$, with $30$ players.

The additive calibration $w_{ijk}=1+\mathbf1_{i=1}+\mathbf1_{j=1}$ with
$m\ge2$ players per sector has $T=m+2$ and core sector shifts
$(t_A,t_B,t_C)\ge0$, $t_A+t_B+t_C=1$, on top of the two exceptional
unit payoffs. All triples lie in optimal perfect matchings and are tight,
which proves this entire-core description. Leximin uses shifts
$(1/3,1/3,1/3)$ and minimax $(0,0,1)$, giving $4m/[3(m+2)]$.
This benchmark is exceeded by Theorem~\ref{thm:main}.

Exact rational checks cover all $256$ binary $2\times2\times2$ arrays,
the sourced empty-core example below, twelve deterministic nonadditive
$3\times3\times3$ arrays, additive cases, and raw recurrence families through
$192$ players. Integral primal matchings and feasible covering allocations
certify the larger examples' matching totals before full staged selector
optimization. Separate implementations check the construction and the
unshifted control. These finite checks support the conventional proof;
they are not the asymptotic argument.

For literature calibration, the published Atay--N\'u\~nez example has layers
$\bigl(\begin{smallmatrix}3&1\\2&5\end{smallmatrix}\bigr)$ and
$\bigl(\begin{smallmatrix}1&4\\5&4\end{smallmatrix}\bigr)$, with unique
diagonal matching total $7$. Lucas's two conditions, as reproduced there,
are $2w_{111}+w_{222}\ge w_{112}+w_{121}+w_{211}$ and
$w_{111}+2w_{222}\ge w_{221}+w_{212}+w_{122}$, conditional on an optimal
diagonal. Here the second reads $11\ge14$, so the core is empty.
The exact calibration reproduces this failure.

Quint initiated the multisided core model; assortative and graph-based
subclasses impose additional restrictions and do not supply a bound for
arbitrary triple arrays. The dated companion source audit found no matching
prior theorem in its bounded search. Full original Quint and Lucas articles
were not available in that audit. Accordingly, this note records an internally
proved structural result, not a certified novelty or field-level breakthrough.

\begin{thebibliography}{9}
\footnotesize\setlength{\itemsep}{2pt}
\bibitem{quint} T.~Quint (1991), The core of an $m$-sided assignment game,
\emph{Games and Economic Behavior} 3, 487--503.
\href{https://doi.org/10.1016/0899-8256(91)90017-9}{doi:10.1016/0899-8256(91)90017-9}.
\bibitem{atay} A.~Atay and M.~N\'u\~nez (2019), A note on the relationship
between the core and stable sets in three-sided markets,
\emph{Mathematical Social Sciences} 98, 10--14.
\href{https://doi.org/10.1016/j.mathsocsci.2018.12.002}{doi:10.1016/j.mathsocsci.2018.12.002}.
\bibitem{lucas} W.~F.~Lucas (1995), Core theory for multiple-sided assignment
games, \emph{Duke Mathematical Journal} 81, 55--65.
\href{https://doi.org/10.1215/S0012-7094-95-08106-X}{doi:10.1215/S0012-7094-95-08106-X}.
\bibitem{assortative} J.~Mart\'inez-de-Alb\'eniz, C.~Rafels and N.~Ybern (2020),
Assortative multisided assignment games: The extreme core points,
\emph{Games and Economic Behavior} 120, 144--153.
\href{https://doi.org/10.1016/j.geb.2020.01.001}{doi:10.1016/j.geb.2020.01.001}.
\end{thebibliography}
\end{document}
